\section{Sealed rounds}\label{app:sealed}

The paper treats games with public moves. What happens when a move can be hidden? Suppose one player chooses a bit and another guesses it. With public moves the guesser wins. Hide the bit, and at any pair of pure actions one of them wishes to switch, so there is no pure equilibrium. Values must therefore retain probabilities.

There is a clean case in which they compose: rounds in which everything is disclosed at the end. The games are now leaves $\at i$ and rounds. A round $\mathsf R_A(\alpha\mapsto G_\alpha)$, with $A=\prod_{i\in I}A_i$ for finite nonempty action sets $A_i$, works as follows. Each role commits privately, all actions are then revealed, and play continues with the public continuation $G_\alpha$. A role with more than one action is \emph{active}: one active role gives public owned choice, and two active roles with Boolean menus give a sealed two-bit round. We can read a round as a multistage game with observed actions, or implement it serially, with the roles committing one at a time and unseen. Let $\Delta$ be the simplex of distributions on $I$, and let $\mathcal E(G)\subseteq\Delta$ be the set of winner distributions of behavioral sequential equilibria (SE) of the serial implementation of $G$; their coordinates are the players' expected utilities. For nonempty sets $P_\alpha\subseteq\Delta$, define
\begin{equation}\label{se:eq:round}
 \mathcal N_A((P_\alpha)_\alpha)=
 \left\{\sum_{\alpha\in A}\left(\prod_i\sigma_i(\alpha_i)\right)p_\alpha:
 \begin{array}{l}
 p_\alpha\in P_\alpha\text{ for every }\alpha,\\
 \sigma\text{ is a mixed Nash equilibrium of the finite game}\\
 \hfill u_i(\alpha)=(p_\alpha)_i
 \end{array}\right\}.
\end{equation}
A value $p_\alpha$ is selected for every tuple, including those of probability zero, and this records the off-path continuation. On public trees mixing adds no possible winner (\cref{se:cor:public}), so the winner-set algebra is exactly right for public games. With hidden moves it is not: strategy alone produces new probabilities (\cref{app:sealed-scalar}), and no finite-valued compositional semantics tracks even who might win (\cref{se:thm:infinite}).

\subsection{The round recursion}

\paragraph{The interface.}
Inactive roles have a singleton action set. No action or signal depending on a commitment becomes available to another role before all actions are bound, and every action tuple, including one with equilibrium probability zero, has its own continuation. The language has no private input types, unrevealed cells crossing a round boundary, common
private random seeds, abort choices, or strategic publication times. Current behavioral randomizations are independent across roles, and the public history can determine the next finite menus and continuation syntax.

Binding and release follow the stable-binding and audience pattern of \cref{sec:disclosure}. For commitment and release orders $i_1,\ldots,i_K$ and $j_1,\ldots,j_K$, a round's deterministic command block is
\[
 \begin{aligned}
 &\Bind_{i_1}(c_{i_1}:A_{i_1};\{i_1\}).\ \cdots\
       \Bind_{i_K}(c_{i_K}:A_{i_K};\{i_K\}).\\
 &\Pub_{j_1}(c_{j_1};I).\ \cdots\ \Pub_{j_K}(c_{j_K};I).
 \end{aligned}
\]
The round language additionally allows randomized binding and tuple-dependent continuation selection, and we assume no translation from the deterministic calculus and no cryptographic or asynchronous realization. Audience-derived views identify exactly the histories with the same preceding public history at each commitment, and no role has a current
private input before its one commitment. Public observations during release are followed by no further choice until the whole tuple is available, so release order has no strategic effect.

As an observed-action multistage game, the language has the standard mixed subgame-perfect semantics \cite[Sections~3.2 and 4.2]{FudenbergTirole1991}. We also consider its serial commitment implementation and use the assessment definition of sequential equilibrium: sequential rationality together with beliefs that are limits of Bayes beliefs of
completely mixed profiles \cite{KrepsWilson1982}.

\begin{proposition}[Round recursion]\label{se:thm:round}
Mixed SPE winner distributions of the observed-action game coincide with SE winner distributions of its serial commitment implementation. For these common values,
\[
 \mathcal E(\at i)=\{\delta_i\},\qquad
 \mathcal E(\mathsf R_A(\alpha\mapsto G_\alpha))=
        \mathcal N_A((\mathcal E(G_\alpha))_\alpha).
\]
\end{proposition}

\begin{proposition}[Properties of round values]\label{se:prop:round-values}
Each value set is nonempty, compact, and semialgebraic over the rationals. The round operators are monotone in their continuation sets and commute with role renaming. Reordering the private commitments within a round preserves the value set.
\end{proposition}

\begin{proof}[Proof of \cref{se:thm:round,se:prop:round-values}]
The standard finite-horizon one-stage-deviation argument gives the
mixed SPE recursion \cite{FudenbergTirole1991}. We prove its equivalence with the serial SE model.
After full publication, each $G_\alpha$ is a proper public subgame. Restricting
an assessment to it gives a sequential equilibrium and a distribution
$p_\alpha$. Any pure change of the current action, followed by the prescribed
continuation, is a legal deviation. The current independent action laws
therefore form a Nash equilibrium for the displayed payoffs.

Conversely choose one continuation equilibrium at every tuple and a
stage Nash profile. Backward induction rules out deviations that also
change later actions: each selected continuation is already optimal
against its opponents, and a current action change cannot improve its
expected continuation payoff. To verify consistency rather than just
subgame optimality, perturb every action law toward its uniform law
with one common positive parameter. At a public round history its reach
factor cancels from Bayes' rule. At a player's commitment, the belief
over preceding hidden commitments is the product of those roles'
perturbed laws. Its limit is the product of the prescribed laws, also
at off-path public histories. Revelation steps offer no choices. This
constructs one consistent assessment for the whole finite game.
The same argument works for every serial order of private commitments.

Nonemptiness follows by choosing continuation distributions and applying
finite mixed Nash existence \cite{Nash1951}. The variables $p_\alpha$ and
$\sigma_i$ range over compact sets. Nash optimality is a finite family
of weak polynomial inequalities, so the feasible set is closed and
its continuous image is compact. The simplex, these inequalities, and
the output equation are rational polynomial conditions. Projection
preserves semialgebraicity, giving the induction. Inclusion of
continuation sets preserves every old feasible choice. Role renaming
only permutes the inequalities and probability coordinates.
\end{proof}

This is the standard backward continuation-payoff construction for games with observed past actions and simultaneous current moves, expressed as an algebra for an explicit disclosure boundary. The continuation $p_\alpha$ remains indexed by its public tuple until the stage incentives have been checked.

\subsection{Public choice}

For a distribution set $P$, put $\operatorname{Win}(P)=\{i:\exists p\in P,\ p_i>0\}$. The one-active-role instance of \eqref{se:eq:round} is
\begin{equation}\label{se:eq:public}
 P\mathbin{\widehat\star_r}Q=
 \left\{\alpha p+(1-\alpha)q:
 \begin{array}{l}
 p\in P,\ q\in Q,\ 0\leq\alpha\leq1,\\
 \alpha>0\Rightarrow p_r\geq q_r,\\
 \alpha<1\Rightarrow q_r\geq p_r
 \end{array}\right\}.
\end{equation}
For singleton inputs it selects the distribution with larger $r$ coordinate, or the full segment between them on a tie.

\begin{corollary}[Public mixed-support conservativity]\label{se:cor:public}
On finite public owned-choice trees,
$\operatorname{Win}(\mathcal E(G))=\den G$. Allowing mixed sequential
equilibria adds no possible winner to the pure SPE observation.
\end{corollary}
\begin{proof}
Pure SPE profiles are sequential equilibria in this perfect-information
language. For the other inclusion induct on trees. If a child has pure
winner set $\{r\}$, its mixed continuation distributions are all
$\delta_r$ by induction. The parent owner $r$ can secure utility one,
so every parent equilibrium also has distribution $\delta_r$.
Otherwise the pure recurrence already includes every possible child
winner; mixing can introduce no others. Nonemptiness of each child's
equilibrium set is used throughout.
\end{proof}

The laws of the winner-set algebra belong to this observation, however, and not to the distributions, which need not even be convex. Let
\[
 T=(\at a\ch c\at b)\ch a\at c.
\]
The left child can give player $a$ any probability $x\in[0,1]$. When $x>0$ the root must choose left, and when $x=0$ the root may mix. Consequently
\begin{equation}\label{se:eq:nonconvex}
 \mathcal E(T)=
 [\delta_a,\delta_b]\ \cup\ [\delta_b,\delta_c].
\end{equation}
In particular $(\delta_a+\delta_c)/2$ is absent. It belongs to $\mathcal E(T\ch b T)$: choose continuation equilibria $\delta_a$ and $\delta_c$ in the two independent copies, and let the indifferent root mix equally. Thus distribution-level public choice is not idempotent, the rewrites of the paper do not carry over to distributions, and convexifying value sets would change the semantics.

\subsection{The scalar fragment}\label{app:sealed-scalar}

With two active roles and every leaf won by one of them, a game is constant-sum. Induction reduces every stage to a finite zero-sum matrix game, whose equilibrium value is unique, so $\mathcal E(G)=\{v(G)\delta_a+(1-v(G))\delta_b\}$, where $v(G)$ is $a$'s winning probability. For public choices owned by $a$ and $b$, the scalar operations are $\max$ and $\min$, and exchanging the roles throughout a component gives $C(x)=1-x$. Consider the round $D(G,H)$:
\begin{center}
\begin{tabular}{c|cc}
 & $b:0$ & $b:1$\\\hline
$a:0$ & $G$ & $\at b$\\
$a:1$ & $\at b$ & $H$
\end{tabular}
\end{center}
The bits are committed privately and both published before entering the selected continuation. Player $a$ wants to match and $b$ to mismatch: if $v(G)=x$ and $v(H)=y$, the stage payoff matrix for $a$ is $\left(\begin{smallmatrix}x&0\\0&y\end{smallmatrix}\right)$. The \emph{scalar fragment} consists of the games built from $\at a$ and $\at b$ by $\ch a$, $\ch b$, $C$, and $D$.

\begin{proposition}[Parallel sum]\label{se:prop:parallel}
The constructor $D$ denotes
\[
 x\parallel y=
 \begin{cases}xy/(x+y),&x+y>0,\\0,&x=y=0.\end{cases}
\]
It is associative, commutative, monotone, and zero-absorbing. On the
unit interval it distributes over both $\min$ and $\max$.
\end{proposition}
\begin{proof}
For $x+y>0$, both players assign probability $y/(x+y)$ to their
action $0$. Against this law, either pure action of the opponent gives
$xy/(x+y)$. These two security guarantees force that value in every
Nash equilibrium. If one diagonal entry is zero the same value follows;
if both are zero every strategy profile has value zero.
For positive arguments, reciprocal values turn parallel sum into
addition. This proves associativity and commutativity; the zero cases
are immediate. Monotonicity follows from the displayed formula, and
a monotone unary map on a chain preserves binary minima and maxima.
\end{proof}

So strategy alone produces new probabilities. The reciprocal-sum value and reciprocal-proportional optimal strategies are the classical diagonal matrix-game formula \cite[Part~II, Section~3.3]{FergusonNotes}, and parallel sum is an established operation in network and matrix theory \cite{AndersonDuffin1969}. The constructor has two Boolean commitments and one occurrence of each continuation, and it introduces neither chance nor numerical rewards at leaves.

\begin{proposition}[Rational values]\label{se:thm:rational}
The scalar values generated by the winner leaves, role exchange, and
$D$ are exactly $\mathbb Q\cap[0,1]$. In fact $C$ and
$U(x)=x\parallel1=x/(1+x)$ suffice with the leaves.
\end{proposition}
\begin{proof}
The operations preserve rationality. For the converse induct on the
denominator of $p/q$, with $0\leq p\leq q$ and $q>0$. Endpoints are
leaves. If $0<p<q$ and $2p\leq q$, write
\[
 p/q=U\bigl(p/(q-p)\bigr).
\]
Otherwise write
\[
 p/q=C\bigl(U((q-p)/p)\bigr).
\]
Each recursive denominator is smaller. Reduction of fractions is not
needed. This gives an explicit synthesis procedure.
\end{proof}

Rational values of finite rational zero-sum games are standard; the point is the small constructor set realizing them with stable Boolean commitments and complete-revelation boundaries.

\begin{corollary}[Decidable scalar equations]\label{se:cor:decidable}
Equations in $0,1,C,\parallel,\min,\max$ hold on rational scalar
values exactly when they hold on the real unit interval. Their validity
is decidable by real quantifier elimination.
\end{corollary}
\begin{proof}
All primitives are continuous. Parallel sum is continuous also at
$(0,0)$ because $0\leq x\parallel y\leq\min(x,y)$.
Rational valuations are dense. For the decision claim, the graph of
parallel sum is given on the unit square by $z(x+y)=xy$, with $z=0$
when $x+y=0$; the other graphs are polynomial or piecewise polynomial.
Term evaluation and a failed equality are therefore expressible in the
first-order theory of real closed fields. Its decision theorem applies.
\end{proof}

\subsection{No finite winner algebra}\label{app:sealed-infinite}

Put $L_0=\at a$ and $L_{n+1}=D(L_n,\at a)$. Then $v(L_n)=1/(n+1)$. For $n\geq0$ let $G_n=L_{n+1}$, so every $G_n$ has possible winners $\{a,b\}$. The winner-set observation cannot tell these games apart, but a third player can. Let $R_m$ be $G_m$ with $c$ in place of $b$, and consider the public context $F_m[-]=[-]\ch a R_m$.

\begin{theorem}[No finite winner algebra]\label{se:thm:infinite}
For all $n,m\geq0$,
\[
 \operatorname{Win}(\mathcal E(F_m[G_n]))=
 \begin{cases}
 \{a,b\},&n<m,\\
 \{a,b,c\},&n=m,\\
 \{a,c\},&n>m.
 \end{cases}
\]
So no finite-valued compositional semantics of the round language can preserve even player $c$'s possible-win observation, already with three roles and Boolean commitments.
\end{theorem}
\begin{proof}
The root owner compares $1/(n+2)$ with $1/(m+2)$, selecting the larger value, or either continuation on equality. This gives the formula. Player $c$ is thus a possible winner exactly when $m\leq n$, so distinct indices respond differently to the contexts $F_m$, and an adequate compositional semantics must assign them distinct values.
\end{proof}

All inputs $G_n$ have the same winner set while the outputs vary, so the possible-winner observation is not compositional on winner sets, and its contextual equivalence has infinite index. This holds with complete revelation after every round and no internal role switches, and extra inactive roles do not affect the construction.
Each step $L_{n+1}=D(L_n,\at a)$ adds one round and only one occurrence of its predecessor, so explicit tree size is linear in $n$. The games $G_n$ have the unique distribution
\[
 p_n=\frac1{n+2}\delta_a+\frac{n+1}{n+2}\delta_b.
\]
The first separator uses rates $1/2$ and $1/3$, against an outside component with probabilities $1/2$ for each of players $a,c$. The first input ties and permits player $b$ to win; the second loses the comparison and excludes player $b$. The obstruction already occurs in the fragment where off-path beliefs and retained private information have been eliminated by
full revelation.

\begin{proposition}[Finite threshold state bound]\label{se:prop:threshold-states}
Let $N\geq0$, and give each of the games $G_0,\ldots,G_N$ of
\cref{se:thm:infinite} a summary, so that games with equal summaries
agree on whether $c\in\operatorname{Win}(\mathcal E(F_m[-]))$ for
$1\leq m\leq N$. Then the $N+1$ summaries are distinct. No
compositionality is assumed.
\end{proposition}
\begin{proof}
By \cref{se:thm:infinite}, $c\in\operatorname{Win}(\mathcal E(F_m[G_n]))$
exactly when $m\leq n$. If $i<j\leq N$, the context $F_{i+1}$ therefore
separates $G_i$ from $G_j$.
\end{proof}

The separation in the infinite family works for every two-player
component, not only for the scalar ones.

\begin{proposition}[Boolean observation of two-player components]\label{se:prop:pair-observation}
Let $a,b,c$ be distinct, and let $G$ and $H$ be games in which only $a$
and $b$ are active and every leaf is won by $a$ or $b$. The following are
equivalent:
\begin{enumerate}
\item $\mathcal E(G)=\mathcal E(H)$, that is, $v(G)=v(H)$;
\item for every context $F$ with Boolean menus,
  $c\in\operatorname{Win}(\mathcal E(F[G]))$ exactly when
  $c\in\operatorname{Win}(\mathcal E(F[H]))$;
\item for every context $F$ with Boolean menus,
  $\operatorname{Win}(\mathcal E(F[G]))=\operatorname{Win}(\mathcal E(F[H]))$.
\end{enumerate}
In the second item the countably many contexts $[-]\ch a R$ suffice,
where $R$ is built from $\at a$ and $\at c$ by $C$ and $U$.
\end{proposition}
\begin{proof}
Here $\mathcal E(G)=\{v(G)\delta_a+(1-v(G))\delta_b\}$, and
by \eqref{se:eq:round} $\mathcal E(F[G])$ depends only on
$\mathcal E(G)$, also with repeated independent occurrences. So the first
item implies the third, for every context, and the third implies the
second. Conversely, let $v(G)<v(H)$ and choose a rational $t$ between
them. By \cref{se:thm:rational}, with $c$ in place of $b$, some $R$ as
above has value $t$, the probability that $a$ wins. In $[-]\ch a R$,
player $a$ selects $R$ against $G$ and $H$ against $R$. As $t<1$, player
$c$ can win after $G$, and it cannot win after $H$. The context has
Boolean menus, since $D$, $C$, and public choice do.
\end{proof}

\begin{corollary}[Scalar full abstraction]\label{se:cor:scalar}
For closed components in the scalar fragment with the same
active pair, equality of scalar values is exactly contextual
equivalence for possible SE winners in the three-role round language.
It suffices to observe the third role's possible victory.
\end{corollary}
\begin{proof}
Equal values give equal singleton distribution sets, which
\eqref{se:eq:round} preserves in every context, including repeated
independent occurrences. Unequal values are separated by
\cref{se:prop:pair-observation}.
\end{proof}

The same contexts recover the order of values. If $v(R)<1$, then $c$
can win in $G\ch a R$ exactly when $v(G)\leq v(R)$: player $a$ selects
$R$ when its value is larger and may select it on a tie. Hence
$v(G)\leq v(H)$ exactly when every such context in which $c$ can win
after $H$ also lets $c$ win after $G$; a rational strictly between
unequal values gives the converse.

\paragraph{What the interface composes.}
The round operator unifies public and sealed choices in one finite binder scheme. Full publication allows continuation equilibria to be selected independently at each public action tuple, and removes a separate belief parameter from the algebra. Before that boundary the commitments cannot be treated as sequential public choices: revealing one matching-pennies
bit before the other's commitment lets the later player dictate the outcome.

The interface summarizes a closed continuation with the fixed winner utilities. Replacing a continuation by another with the same $\mathcal E$ is sound. Replacing the terminal rewards inside a previously summarized component requires a continuation-dependent response semantics. A pending secret, private type, selective audience, or strategic release leaves an
information state crossing the interface, and marginal winner probabilities then do not specify the needed joint conditional laws and consistent beliefs.
